
% \subsection{Early Methods for Analyzing RCS Data}



Among the earliest methods designed for RCS data was that proposed by Scott and Smith~\cite{scott1974repeatedsurveys}, in which dependent variables are aggregated at the time step level and modeled as a time series with added noise:


\begin{equation}
    \label{eq:scott}
    y_t = \theta_t + e_t, \qquad \text{Var}(e_t) = s_t^2
\end{equation}

In this model, surveys are conducted at regular time intervals and $\theta_t$ represents a population mean changing over time. The analytical method proposed therein introduces different possibilities for the covariance structures of $\{\theta_t\}$ and $\{e_t\}$, but is ultimately concerned with the use of all previously observed aggregates $\{y_t\}_{t=1}^T$ to estimate $\theta_T$. As our objective is to improve the estimation and of covariate effects and their standard errors, we must look to more highly parameterized models.

Such a model was developed by DiPrete and Grusky \cite{diprete1990multilevel}, in which added complexity and flexibility is added to a standard fixed effects model by having the effects of individual-level variables depend on macro-level variables:

\begin{equation}
    \label{eq:diprete}
    y_{ti} = X_{ti}\beta_t + \varepsilon_{ti}, \qquad \beta_t = \mathbf{Z}_t \mathbf{\gamma}, \qquad \varepsilon_{ti} \sim \text{WN}(0,\sigma^2)
\end{equation}

Such a model requires the presence of macro-level variables that vary over time and explain changing effects of individual-level variables. While this formulation has a broad range of applications, it is only appropriate in cases where temporal correlation can be assumed to be a function of known time-varying predictors, rather than an independent random process.
